\chapter{Introduction}
\label{ch:introduction}

Autonomous landing of unmanned aerial vehicles (UAVs) remains a safety-critical challenge
in aerospace engineering.
Unlike take-off, landing requires the vehicle to bring all translational velocities to zero
simultaneously with altitude, leaving little margin for error in the terminal phase.
Classical control approaches rely on precise sensor feedback and pre-programmed deceleration
profiles \cite{flores2012}; however, biological fliers achieve reliable landing without
explicit velocity sensors, using solely the optic flow available in the visual field.

Honeybees (\emph{Apis mellifera}), hoverflies, and dragonflies are known to regulate their
approach by keeping the rate of change of the \emph{tau} variable
($\tau = z / |v_z|$, the optical time-to-contact) approximately constant at
$\dot{\tau} \approx -0.5$ \cite{lee1976, wagner1982}.
If this condition holds throughout descent, velocity decreases proportionally with altitude,
guaranteeing a soft touchdown regardless of initial conditions \cite{lee1976}.
This makes tau-regulation a robust, sensor-minimal, and provably safe landing
strategy---properties highly desirable for autonomous UAVs.

In parallel, deep reinforcement learning (RL) has demonstrated the ability to discover
complex control strategies through environmental interaction alone \cite{schulman2017ppo}.
A natural question is whether augmenting an RL agent with the biological percept $\tau$
and a corresponding reward signal will (i)~improve landing quality and
(ii)~induce tau-regulation behaviour similar to that observed in insects.

This report investigates this question through the following contributions:
\begin{enumerate}
  \item A custom 2-D drone landing environment (\texttt{DroneEnvTau2D}) extending
        OpenAI Gymnasium \cite{towers2024} with a tau observation and bio-inspired
        shaping reward based on Lee's theory \cite{lee1976}.
  \item A multi-seed experimental protocol quantifying the improvement in landing
        performance and bio-inspired behaviour with statistical rigour (Wilcoxon
        signed-rank test, bootstrap 95\% confidence intervals).
  \item A zero-shot wind robustness evaluation showing that tau shaping preserves
        landing quality under atmospheric disturbances the agent was never trained on.
  \item A sensitivity analysis identifying the optimal reward shaping coefficient
        governing the tau-regulation/task-success trade-off.
\end{enumerate}

The report is structured as follows.
Chapter~\ref{ch:theory} introduces the theoretical background.
Chapter~\ref{ch:method} describes the environment, agent, and training procedure.
Chapter~\ref{ch:experiments} defines the experimental design.
Chapter~\ref{ch:results} presents quantitative results.
Chapter~\ref{ch:analysis} analyses the learned strategy.
Chapter~\ref{ch:conclusion} concludes and discusses future work.
The source code is publicly available at \url{https://github.com/stosvanlynden/insect-landing}.
